\documentclass[twocolumn]{aastex631}

\usepackage{amsmath}
\usepackage{multirow}
\usepackage{natbib}
\usepackage{graphicx} 
\usepackage{aas_macros}

\begin{document}

The quantum stability of Horndeski theories, particularly ensuring ghost-freedom beyond classical approximations, remains a formidable challenge due to their complex structure. This paper introduces a novel approach by identifying a non-trivial sub-sector of Horndeski theories that possesses a hidden disformal symmetry, allowing it to be mapped to a simpler, demonstrably stable canonical scalar field minimally coupled to General Relativity. Through systematic analytical calculations, we explicitly derive the functional forms of the Horndeski functions for this sub-sector, defined by a specific disformal metric transformation. A one-loop perturbative analysis around a Friedmann-Robertson-Walker background in the equivalent, simpler frame immediately confirms the absence of ghost and tachyonic instabilities, properties directly inherited by the original Horndeski model. Crucially, this disformally protected sub-sector automatically satisfies the stringent observational constraint of gravitational waves propagating at the speed of light ($c_T^2=1$) and predicts distinct, testable signatures in the growth of large-scale structure. Furthermore, we argue that this disformal equivalence, being an invertible field redefinition, grants non-perturbative quantum protection, ensuring the absence of ghosts to all loop orders by inheriting the fundamental unitarity of the canonical target theory. This framework thus establishes a theoretically robust and observationally viable class of modified gravity theories, offering a concrete realization of mechanisms for radiative stability.

\end{document}
                